% =====================================================================
% SPARCS manuscript
% Compile with pdfLaTeX. Overleaf: upload this file plus (optionally)
% table1.tex and the figure PNGs written by figures.py.
%
% Table 1 lives in table1.tex and is regenerated by:
%     python experiments.py ...
%     python figures.py
% Figures are pulled in only if the PNG exists; otherwise a visible
% placeholder box is typeset so a missing figure cannot slip through.
% =====================================================================
\documentclass[11pt]{article}

\usepackage[margin=1in]{geometry}
\usepackage{amsmath,amssymb}
\usepackage{graphicx}
\usepackage{booktabs}
\usepackage{caption}
\usepackage{enumitem}
\usepackage{algorithm}
\usepackage{algpseudocode}
\usepackage{xcolor}
\usepackage{float}
\usepackage[numbers,sort&compress]{natbib}
\usepackage[hidelinks]{hyperref}

\captionsetup{font=small,labelfont=bf}
\setlist{nosep,leftmargin=*}

% Include a figure if it has been generated, otherwise leave an obvious gap.
\newcommand{\resultfigure}[3]{%
  \begin{figure}[H]\centering
  \IfFileExists{#1}{\includegraphics[width=#2\textwidth]{#1}}{%
    \fbox{\parbox[c][3.4cm][c]{0.85\textwidth}{\centering
      \texttt{\detokenize{#1}} not found.\\[3pt]\small
      Run \texttt{python experiments.py} then \texttt{python figures.py}
      to generate this figure.}}}
  \caption{#3}\end{figure}}

\title{\vspace{-2em}A Comparative Framework for Robust Variable Selection:\\
Integrating Stability, Rank-Based Screening, and Adaptive Regularization}
\author{Kashfi Rehan\\\small The Bronx High School of Science}
\date{}

\begin{document}
\maketitle

\begin{abstract}
\noindent High-dimensional biomedical datasets, in which features vastly outnumber samples, pose a
major challenge for precision medicine: identifying truly informative molecular features among
thousands of candidates. Feature selection methods span diverse paradigms, including adaptive
regularization, iterative screening, and stability selection, but practical performance across
varying data complexity remains unclear. We conduct an empirical comparison of four methods on
synthetic datasets with known ground truth and real cancer genomics data, with particular focus on
SPARCS (Stable Permutative Adaptive Rank-Based Correlation Screening), the novel integrated
framework and contribution of this paper, combining rank-based screening, nonlinear dependence
detection via the Randomized Dependence Coefficient, permutation testing, and batch Complementary
Pairs Stability Selection. We evaluated methods on two synthetic regression datasets differing in
nonlinearity and complexity and two cancer gene expression datasets from the CuMiDa repository. In
moderately complex synthetic data, SPARCS achieved the highest feature recovery while maintaining
competitive predictive performance. Under extreme nonlinearity and noise, however, adaptive Elastic
Net outperformed all screening-based methods, improving predictive accuracy as complexity increased,
highlighting the robustness of strong $L_1+L_2$ regularization against nonlinearity especially for
multiclass classification. All methods achieved comparable performance on binary classification but
differed substantially in multiclass performance and computational efficiency. Feature set overlap
across methods was low, indicating divergent but valid selection priorities. SPARCS demonstrated
adaptive sparsity, selecting a substantially smaller feature set for multiclass classification while
achieving competitive accuracy. These results provide practical guidance for method selection in
high-dimensional biomedical studies and highlight fundamental trade-offs between sensitivity,
stability, computational efficiency, and robustness to complexity.
\end{abstract}

\section{Introduction}

Modern biomedical research compiles and generates datasets where the number of measured variables
(features) often vastly exceeds the number of samples, a scenario known as the ``high-dimensional''
or ``large $p$, small $n$'' problem \citep{hastie2009,saeys2007}. For example, a gene expression
study might measure 20{,}000 genes across only a few hundred patients \citep{saeys2007}. In such
settings, most measured features are irrelevant noise, and identifying the truly informative subset
is both statistically challenging and computationally demanding
\citep{hastie2009,guyon2003}. Selecting too many features leads to overfitting, poor generalization,
and unreliable interpretability, while missing important features results in lost predictive power
and scientific insight \citep{guyon2003}.

Traditional statistical methods are not applicable in high dimensions: classical regression assumes
more samples than features, and standard hypothesis testing suffers from multiple testing problems
when screening thousands of variables simultaneously \citep{buhlmann2011}. Moreover, biomedical data
in particular poses unique challenges. Features are often correlated due to biological mechanisms
such as co-expression of genes in regulatory pathways \citep{barabasi2011}, feature-outcome
relationships may be nonlinear or involve complex interactions \citep{zhao2014,liu2009}, and data
distributions are usually heavy-tailed with outliers due to measurement noise and biological
heterogeneity \citep{anders2010}.

\subsection{Existing Approaches and Their Limitations}

Feature selection literature has developed various families of methods to address these challenges,
each differing in theoretical foundations and practical trade-offs
\citep{guyon2003,chandrashekar2014,li2017}.

\paragraph{Regularization-based methods.} These perform feature selection by adding penalty terms to
regression coefficients during model fitting, setting irrelevant features to exactly zero. The LASSO
(Least Absolute Shrinkage and Selection Operator) uses an $L_1$ penalty that pushes strong sparsity
\citep{tibshirani1996}. The Elastic Net extends this by combining $L_1$ and $L_2$ (Ridge) penalties
to handle correlated features \citep{zou2005}. Adaptive variants incorporate per-feature penalties
based on initial estimates, improving selection consistency and demonstrating oracle properties under
specific conditions \citep{zou2006}. These methods are computationally efficient and theoretically
backed, but they assume linearity and can be sensitive to the chosen regularization strength
\citep{buhlmann2011}. Thus, they struggle with nonlinear relationships and may miss important
features when correlations are strong.

\paragraph{Screening-based methods.} These take a different approach by first vastly reducing the
feature space through marginal correlations before applying more sophisticated selection algorithms.
Sure Independence Screening (SIS) and its iterative variant (ISIS) use correlations to quickly
eliminate most irrelevant features, then iteratively refine the selection
\citep{fan2008,fan2009,fan2010}. These methods are extremely fast and provide theoretical guarantees
under certain conditions. For one, they reduce computational complexity from $O(p)$ model fits to a
single-pass marginal screening step. However, standard implementations use Pearson correlation, which
is sensitive to outliers and assumes linear, monotonic relationships. Robust variants have been built
to address these limitations, a notable one involving Spearman rank correlations, which can resist
outliers and capture monotonic nonlinear relationships \citep{li2012,wang2013,croux2010}. Even with
these improvements, correlation-based screening measures only pairwise linear or monotonic
associations, often missing features with more complex nonlinear or interaction effects.

\paragraph{Stability-based methods.} These address a different concern: selection stability and false
discovery control \citep{meinshausen2010,nogueira2018}. Even if a method illustrates good predictive
performance, the selected feature subset may change dramatically with small data perturbations. This
is important for scientific interpretability and clinical translation: a signature that varies wildly
across patient cohorts lacks biological meaning and practical utility
\citep{nogueira2018,eindor2005,michiels2005}. Stability selection, introduced by Meinshausen and
B\"uhlmann, repeatedly applies a base selector to random dataset subsamples and only retains features
selected consistently across subsamples according to a chosen stability threshold
\citep{meinshausen2010}. Complementary Pairs Stability Selection (CPSS), developed by Shah and
Samworth, enhances this process with finite-sample error control bounds \citep{shah2013}. This
framework provides rigorous statistical control that most other methods lack. However, stability
selection is computationally expensive, requiring many repeated applications of the base selector,
and tends to be conservative, often selecting fewer features than optimal for prediction
\citep{hofner2015}.

\subsection{Integration and Nonlinear Dependence}

More recent methodological developments have been exploring ways to integrate the strengths of
different tools while balancing their individual limitations. One particularly promising direction
involves incorporating measures of nonlinear dependence into the screening process. The Randomized
Dependence Coefficient (RDC), introduced by Lopez-Paz and colleagues, provides a nonparametric
measure that captures nonlinear, non-monotonic relationships missed by correlation
\citep{lopezpaz2013}. RDC works by transforming features to copula space, applying random projections
through trigonometric feature maps, and subsequently computing the canonical correlation of those
projections. Importantly, RDC can detect relationships like quadratic effects, threshold phenomena,
and complex interaction patterns that standard correlation entirely misses \citep{lopezpaz2013}.

Integrating RDC into iterative screening frameworks opens doors to combining computational efficiency
with successfully capturing complex biological relationships. Pairing RDC with permutation-based
statistical calibration can provide nonparametric $p$-values without distributional assumptions
\citep{ernst2004}. When further combined with batch stability selection \citep{hofner2015} and
adaptive threshold mechanisms \citep{salvador2004}, such hybrid approaches offer a promising path
towards robust, powerful, and computationally practical feature selection for high-dimensional
biomedical data.

\subsection{Empirical Validation and Benchmarking}

Despite extensive theoretical work on feature selection methods, empirical comparisons on realistic
biomedical data remain surprisingly limited and often incomplete
\citep{saeys2007,chandrashekar2014,michiels2005}. Many studies evaluate methods on only one or two
datasets, focus exclusively on predictive performance while ignoring selection stability, or fail to
compare across different methodological families. This is especially problematic in cancer genomics,
where the heterogeneity of biological mechanisms and data collection protocols means that method
performance can vary substantially across applications
\citep{saeys2007,salvador2004,feltes2019}.

To address these challenges, rigorous benchmarking studies require both synthetic datasets with known
ground truth (enabling direct assessment of feature recovery) and real datasets reflecting practical
applications (assessing real-world performance). The Curated Microarray Database (CuMiDa) repository
provides an excellent resource for such benchmarking \citep{feltes2019}. Feltes and colleagues
manually curated 78 cancer microarray datasets from over 30{,}000 studies in the Gene Expression
Omnibus, applying rigorous filtering, background correction, normalization, and quality control. Each
dataset was carefully processed to remove technical artifacts while preserving biological signals,
making them ideal for machine learning algorithm comparison \citep{feltes2019}.

\subsection{Research Question and Contributions}

This study addresses the following question: how do modern feature selection paradigms, namely
adaptive regularization, iterative rank-based screening, stability selection, and integrated hybrid
approaches, compare in their ability to identify truly informative features and achieve strong
predictive performance on high-dimensional biomedical data?

To answer this question, we implemented and systematically compared four distinct feature selection
methods representing distinct paradigms, the first of which is a novel contribution of this paper.

\begin{enumerate}
\item \textbf{SPARCS} (Stable Permutative Adaptive Rank-Based Correlation Screening). A novel
integrated framework combining rank-based initial screening for robustness, RDC-based nonlinear
dependence detection, permutation-based statistical testing for rigorous calibration, batch CPSS for
stability without excessive computation, and adaptive mechanisms for efficiency.

\item \textbf{Iterative Sure Independence Screening with Rank-Based Spearman Correlations}
(Spearman-ISIS). Iterative greedy screening methodology using robust rank correlations for marginal
screening and conditional selection \citep{fan2009,li2012,wang2013}. This serves as a baseline for
modern screening approaches.

\item \textbf{Complementary Pairs Stability Selection} (CPSS-Spearman). Pure stability selection
methodology, wrapping Spearman-ISIS with Complementary Pairs Stability Selection
\citep{meinshausen2010,shah2013}. This provides conservative selections with formal error bounds.

\item \textbf{Adaptive Stochastic Gradient Descent Elastic Net} (ASGDR for regression, ASGDC for
classification). Efficient adaptive regularization using feature-specific penalties based on initial
estimates \citep{zou2005,zou2006}. This combines the oracle properties of adaptive methods with
computational efficiency through SGD optimization.
\end{enumerate}

Each method was evaluated using a consistent framework: all methods perform feature selection on
training data, then the selected features are used to train a LightGBM model with hyperparameter
tuning \citep{ke2017}. This ensures fair comparison, as all methods are judged by the same powerful
downstream predictor, isolating the quality of feature selection itself.

\section{Methods and Framework}

\resultfigure{figures/fig0_pipeline.png}{0.95}{The comparative feature selection pipeline. Each
method receives the same training split, and every selected subset is passed to the same tuned
LightGBM predictor.}

\subsection{Problem Formulation and Notation}

Consider a dataset with $n$ samples and $p$ features, where $p \gg n$ (the high-dimensional setting).
Let $X \in \mathbb{R}^{n \times p}$ denote the feature matrix, where each row represents a sample
(e.g., a patient) and each column represents a feature (e.g., a gene's expression). Let
$\mathbf{y} \in \mathbb{R}^{n}$ denote the response vector (continuous for regression, categorical
for classification).

The feature selection problem is to identify a subset $S \subseteq \{1, 2, \dots, p\}$ of relevant
features such that $S$ contains all (or realistically most) truly informative features (high
sensitivity/recall); $S$ excludes all or most noise features (high specificity/precision); and a
predictive model trained on $X_S$ achieves strong generalization on test data.

Let $S^\star$ denote the true set of informative features. For synthetic data where $S^\star$ is
known, we evaluate
\begin{equation}
\mathrm{TPR} = \frac{|S \cap S^\star|}{|S^\star|}, \qquad
\mathrm{FDR} = \frac{|S \setminus S^\star|}{|S|}, \qquad
\mathrm{Precision} = \frac{|S \cap S^\star|}{|S|}.
\end{equation}
For real data (unknown $S^\star$), we evaluate indirectly through predictive performance on held-out
test sets.

\subsection{Feature Selection Methods}

\subsubsection{SPARCS}

SPARCS is a novel integrated framework combining multiple robust techniques meant for the
high-dimensional, low-sample-size biomedical data setting. It integrates robust rank-based marginal
screening to account for outliers and heavy-tailed distributions, nonlinear dependence detection to
capture complex feature-outcome relationships, permutation-based statistical calibration to control
false discoveries, stability selection to balance sensitivity, specificity, and computational
efficiency, and adaptive pool shrinking. Throughout, the internal working model is the Adaptive
Elastic Net of Section~\ref{sec:asgd}: it fits the current selected set, supplies the residual that
drives the next screening round, and provides the coefficients that stability selection counts. The
algorithm proceeds as follows.

\paragraph{Step 1: Rank-based Sure Independence Screening warm start.} For each feature $j$,
rank-transform $X_{:,j}$ and $\mathbf{y}$ and compute the Spearman correlation
$\rho_j = \mathrm{corr}(R_{X[:,j]}, R_{\mathbf{y}})$. The working set is initialized with the top $d$
features by $|\rho_j|$, where $d = \lceil n / \log n \rceil$ following the standard SIS rate
\citep{fan2008}. This warm start supplies the first residual; the iterative screening in Steps 2--8
continues to operate over all $p - |S_t|$ unselected features rather than being confined to the
initial $d$.

\paragraph{Step 2: Iterative residual-based refinement.} At iteration $t$ with current selected set
$S_t$, fit the Adaptive Elastic Net on $X_{S_t}$ to obtain predictions $\hat{\mathbf{y}}$ and form
the working residual $\mathbf{r}$. For regression $\mathbf{r} = \mathbf{y} - \hat{\mathbf{y}}$; for
binary classification $\mathbf{r} = \mathbf{y} - \hat{P}(y=1 \mid X)$; for multiclass classification
$r_i = 1 - \hat{P}(y = y_i \mid x_i)$, the residual probability mass on the true class.

\paragraph{Step 3: Staged RDC screening.} Computing the Randomized Dependence Coefficient against
every one of the $p - |S_t|$ remaining features at every iteration dominates the cost of the method,
so screening is staged. Stage~1 ranks the remaining features by $|\rho_j|$ against $\mathbf{r}$ and
retains the top $M_{\mathrm{pre}}$ (default 500). Stage~2 computes
\begin{equation}
\mathrm{RDC}(X_{:,j}, \mathbf{r}) \in [0, 1]
\end{equation}
for each survivor and retains the top $M_{\mathrm{rdc}}$ (default 200, shrunk per Step~8). RDC
transforms both marginals to the empirical copula, applies $k$ random sine and cosine projections
with bandwidth $s$, whitens each projected space, and returns the largest canonical correlation
between them \citep{lopezpaz2013}. We use $k = 20$ and $s = 1$, which keeps the statistic near zero
under independence while recovering quadratic, saturating, and threshold dependence that rank
correlation misses. The whitening step is what makes the returned quantity a correlation on $[0,1]$
rather than a scale-dependent covariance norm.

\paragraph{Step 4: Permutation calibration.} For the $M_{\mathrm{rdc}}$ surviving candidates, permute
the residual $n_{\mathrm{perm}}$ times (default 30), recompute RDC on each permuted replicate, and
form the empirical $p$-value
\begin{equation}
p_{\mathrm{emp}} = \frac{1 + \#\{\text{permuted RDC} \ge \text{observed RDC}\}}{1 + n_{\mathrm{perm}}}.
\end{equation}
For classification the permutation is stratified within class, so the null preserves the class
structure and only breaks the feature-residual association. Note that the attainable minimum is
$1/(1 + n_{\mathrm{perm}})$, so $n_{\mathrm{perm}}$ must satisfy
$1/(1 + n_{\mathrm{perm}}) \le \alpha$; at $\alpha = 0.05$ this requires
$n_{\mathrm{perm}} \ge 19$, which the default of 30 satisfies. Candidates with
$p_{\mathrm{emp}} \le \alpha$ are retained. A hybrid rule additionally retains the
$m_{\min}$ (default 10) strongest candidates by RDC regardless of significance, which guarantees the
iteration has something to score when the permutation null is not separable at the chosen level; the
stability and magnitude filters of Steps 6 and 7 remain in force, so this floor widens the candidate
pool without weakening the admission criteria.

\paragraph{Step 5: Batch CPSS stability.} For $B$ complementary subsample pairs (default 30),
randomly split the data into complementary halves of size $n/2$ and fit the working model once per
half on the full candidate set. A candidate is credited when its coefficient magnitude exceeds
$10^{-8}$. The stability score of feature $j$ is the proportion of the $2B$ half-samples that select
it. Because every candidate is evaluated inside a single fit rather than one fit per candidate, the
batch scheme costs $O(2B)$ model fits instead of $O(B \cdot |\text{candidates}|)$.

\paragraph{Step 6: Adaptive stability threshold.} Sort the stability scores in decreasing order,
lightly smooth the profile, and take the second difference to locate the point of maximum curvature
\citep{salvador2004}. The threshold is set at that elbow and floored at $\tau = 0.6$, so the
data-driven threshold can only ever be more conservative than the fixed one.

\paragraph{Step 7: Multi-add.} Select the top $k_{\mathrm{add}}$ stable candidates (default 3) that
clear both the adaptive stability threshold and a minimum RDC magnitude. The RDC threshold is 0.2 for
synthetic regression data and 0.1 for the real classification datasets.

\paragraph{Step 8: Adaptive pool shrinkage.} The cap on the candidate pool shrinks geometrically,
\begin{equation}
K_{\mathrm{cp}}(t) = \max\!\left(K_{\min},\ \left\lfloor K_{\mathrm{init}} \cdot \gamma^{\,t-1} \right\rfloor\right),
\end{equation}
with shrinkage rate $\gamma = 0.9$ and floor $K_{\min} = 50$. Later iterations screen a narrower pool,
on the reasoning that the strongest residual associations surface early.

\paragraph{Step 9: Stopping criteria.} Iteration halts when (i) no candidate clears both thresholds,
(ii) the feature budget $|S| = $ \texttt{max\_features} is reached, or (iii) the candidate pool is
exhausted. Criterion (i) is what produces data-driven sparsity: SPARCS stops on its own rather than
padding the set out to a fixed target.

\resultfigure{figures/fig0_sparcs.png}{0.8}{The complete SPARCS algorithm. Steps 2--8 repeat until
one of the three stopping criteria fires.}

\subsubsection{Spearman-ISIS}

Spearman-ISIS adapts Iterative Sure Independence Screening to use Spearman rank correlations for
robustness to outliers and monotonic nonlinearities \citep{fan2008,fan2009,fan2010}.

\begin{enumerate}
\item Initialize $S_0 = \emptyset$ and $\mathbf{r}_0 = \mathbf{y} - \bar{\mathbf{y}}$.
\item For $t = 1, 2, \dots$:
  \begin{enumerate}
  \item For each $j \notin S_t$, compute $\rho_j = \mathrm{spearman}(X_{:,j}, \mathbf{r}_{t-1})$.
  \item Select the top $d = \lceil n / \log n \rceil$ candidates by $|\rho_j|$.
  \item Fit the Adaptive Elastic Net on $S_{t-1} \cup \{\text{candidates}\}$.
  \item Greedily commit the single candidate with the largest coefficient magnitude to form $S_t$.
  \item Refit on $S_t$ and update the residual $\mathbf{r}_t$.
  \item Stop when the largest candidate magnitude falls below $10^{-6}$, the feature budget is
  reached, or all features have been considered.
  \end{enumerate}
\end{enumerate}

Key characteristics: iterative residual-based screening, greedy one-per-iteration selection,
rank-based robustness, and computational efficiency.

\subsubsection{CPSS-Spearman}

CPSS-Spearman applies pure Complementary Pairs Stability Selection using Spearman-ISIS as the base
selector \citep{meinshausen2010,shah2013}.

\begin{enumerate}
\item For $b = 1, \dots, B$ pairs (default $B = 30$):
  \begin{enumerate}
  \item Randomly partition the data into disjoint halves $A_b$ and $B_b$ of size $n/2$.
  \item Apply the base selector to each half to obtain
  $S_{A_b} = \text{Spearman-ISIS}(X_{A_b,:}, \mathbf{y}_{A_b})$ and
  $S_{B_b} = \text{Spearman-ISIS}(X_{B_b,:}, \mathbf{y}_{B_b})$.
  \item Feature $j$ is counted in pair $b$ if $j \in S_{A_b}$ or $j \in S_{B_b}$.
  \end{enumerate}
\item Aggregate selection frequencies
$\pi_j = \tfrac{1}{B}\sum_{b=1}^{B} \mathbb{I}\{j \text{ selected in pair } b\}$.
\item Return $S = \{j : \pi_j \ge \tau\}$, with stability threshold $\tau = 0.6$ by default.
\end{enumerate}

\paragraph{Error control.} Let $\hat{q}$ be the average number of features selected per
half-subsample. For $\tau > 0.5$ \citep{shah2013},
\begin{equation}
\mathbb{E}[\#\text{false positives in } S] \le \frac{\hat{q}^{\,2}}{(2\tau - 1)\,p},
\end{equation}
a conservative, distribution-free, finite-sample bound requiring no asymptotic assumptions. The
accompanying implementation reports $\hat{q}$ and the resulting bound alongside every CPSS run.

\subsubsection{ASGD (Adaptive Elastic Net)}
\label{sec:asgd}

ASGD implements the Adaptive Elastic Net trained efficiently via stochastic gradient descent
\citep{tibshirani1996,zou2005,zou2006}, with objective
\begin{equation}
\min_{\beta} \ \frac{1}{2n}\|\mathbf{y} - X\beta\|_2^2
+ \lambda\!\left[(1-\alpha)\|\beta\|_2^2 + \alpha \sum_j w_j |\beta_j|\right],
\qquad
w_j = \frac{1}{\left(\left|\hat{\beta}_j^{\mathrm{init}}\right| + \epsilon\right)^{\gamma}},
\end{equation}
where $\hat{\beta}^{\mathrm{init}}$ are initial ridge estimates, $\gamma = 1$ is the adaptive
exponent, and $\epsilon = 10^{-8}$. Larger initial estimates receive smaller penalties, improving
selection consistency with oracle properties.

\paragraph{Implementation through rescaling.} The weighted penalty is removed by a change of
variables, which lets a standard fast Elastic Net solver do the work.
\begin{enumerate}
\item Fit the initial model to obtain $\hat{\beta}^{\mathrm{init}}$ and compute the weights $w_j$.
For regression this is a ridge fit with $\alpha = 0.1$; for classification it is an
$L_2$-penalized logistic model with the same strength, and for multiclass problems the per-feature
initial magnitude is the norm of the coefficient vector across classes.
\item Compute rescaling factors $s_j = 1/w_j$.
\item Rescale the feature matrix, $\tilde{X}_{:,j} = X_{:,j} / w_j$.
\item Fit a standard Elastic Net on $\tilde{X}$ using stochastic gradient descent.
\item Recover coefficients $\beta_j = \tilde{\beta}_j / w_j$.
\end{enumerate}
This transformation is exact: $X\beta = \tilde{X}\tilde{\beta}$ leaves the fitted values unchanged,
and $\sum_j w_j |\beta_j| = \sum_j |\tilde{\beta}_j|$ turns the weighted $L_1$ term into a plain
$L_1$ term, so the rescaled problem a standard solver sees is the adaptive problem we intended to
solve. For feature selection, features are ranked by adaptive importance,
\begin{equation}
\mathrm{Importance}_j = \frac{|\beta_j|}{w_j},
\end{equation}
and the top $k$ are returned, with $k$ set to the same feature budget given to the other methods.

\subsection{Datasets}

\subsubsection{Synthetic Regression Datasets}

Two synthetic datasets (s1, s2) with continuous targets were generated. Each contains 10{,}000
gene-like features, 500 patient-like samples, and 80 truly informative features providing known
ground truth. Features are generated in equicorrelated blocks of 50 with within-block correlation
$\rho = 0.3$ by construction (empirically $0.265$ in the realized draw), mimicking gene modules.
Feature distributions are mixed: 30\% of columns are pushed to heavy-tailed ($t_3$) or right-skewed
(log-normal) marginals through the probability integral transform, which changes the marginals while
preserving the block dependence structure, reflecting heterogeneous gene expression behavior.

s1 makes 20\% of the informative features act through a nonlinear link and s2 makes 80\%, cycling
through quadratic, logarithmic, saturating (tanh), and threshold effects. s1 contains 5 pairwise
interaction terms $X_i \times X_j$ drawn from the informative set, while s2 contains 10. Noise is
Gaussian with standard deviation set to a fixed multiple of the signal standard deviation: 0.35 for
s1, giving a realized signal-to-noise ratio of 8.48, and 0.50 for s2, giving 4.42. Each dataset was
split 80/20 into train and test sets, stratified by response quantiles, to permit both
feature-recovery evaluation on the training split and held-out predictive evaluation on the test
split. The known ground truth enables direct assessment of feature recovery through TPR and FDR.

\subsubsection{Real Curated Cancer Datasets from CuMiDa}

\textbf{Leukemia\_GSE28497} is an imbalanced multiclass classification dataset whose target contains
7 leukemia subtypes (B-CELL\_ALL, B-CELL\_ALL\_TCF3-PBX1, B-CELL\_ALL\_HYPERDIP,
B-CELL\_ALL\_HYPO, B-CELL\_ALL\_MLL, B-CELL\_ALL\_T-ALL, B-CELL\_ALL\_ETV6-RUNX1), with 281 samples
and 22{,}284 gene probes. It requires identification of subtle gene expression differences across
imbalanced subtypes.

\textbf{Liver\_GSE14520\_U133A} is a binary classification dataset (hepatocellular carcinoma versus
normal) with 357 samples and 22{,}278 gene probes.

\paragraph{Preprocessing.} Both real datasets underwent duplicate-probe removal, median imputation of
any missing values, removal of zero-variance features, StandardScaler normalization (zero mean, unit
variance) fitted on the training split only, and a 70/30 train/test split with stratification. CuMiDa
datasets were already extensively curated with background correction, normalization, and quality
control \citep{feltes2019}.

\subsection{Evaluation Framework}

\subsubsection{Two-Stage Selection and Prediction Pipeline}

All methods were evaluated using a consistent two-stage pipeline that separates feature selection
from downstream prediction.

\textbf{Stage 1 (feature selection).} Each method $M$ receives $(X_{\mathrm{train}},
\mathbf{y}_{\mathrm{train}})$ and outputs a selected subset $S_M$. Methods use recommended literature
parameters to avoid overfitting the selection step to the datasets.

\textbf{Stage 2 (prediction evaluation).} A LightGBM model is trained on
$X_{\mathrm{train}}[:, S_M]$ with hyperparameter tuning under 5-fold cross-validation, then evaluated
on the held-out test set $X_{\mathrm{test}}[:, S_M]$. The search space is
\begin{itemize}
\item Tree structure: \texttt{n\_estimators} $\in \{50, 100, 200\}$, \texttt{max\_depth} $\in
\{3, 5, 7, 10, -1\}$, \texttt{num\_leaves} $\in \{15, 31, 63, 127\}$
\item Learning dynamics: \texttt{learning\_rate} $\in \{0.01, 0.1, 0.2, 0.5, 0.7, 1\}$,
\texttt{min\_child\_samples} $\in \{10, 20, 30, 50\}$
\item Regularization: \texttt{subsample} $\in \{0.6, \dots, 1.0\}$, \texttt{colsample\_bytree} $\in
\{0.6, \dots, 1.0\}$, \texttt{reg\_alpha} $\in \{0.01, 0.1, 0.2, 0.5, 0.7, 1.0\}$,
\texttt{reg\_lambda} $\in \{0, 0.01, 0.1, 0.2, 0.5, 0.7, 1.0\}$
\end{itemize}
Exhaustive enumeration of this space is roughly $1.5 \times 10^{6}$ configurations, which at 5 folds
would require on the order of $10^{7}$ model fits per method and is not tractable. We therefore use
randomized search over the same space with 60 sampled configurations per method, drawn under a fixed
seed so that every method sees the identical budget and the identical candidate stream. Optimization
targets $R^2$ for regression, AUC-ROC for binary classification, and macro F1 for multiclass.

All methods are judged by the same downstream predictor, isolating feature selection quality.
LightGBM is chosen for its performance on tabular data, computational efficiency, and ability to
handle nonlinear relationships \citep{ke2017}.

\subsubsection{Replication and Uncertainty}
\label{sec:replication}

Every configuration is replicated across five seeds. A seed controls the
train/test split and the internal randomization of every selector, so a
replicate is a full end-to-end repetition rather than a re-initialization of one
component. The underlying data is held fixed across replicates, which keeps the
same feature universe on offer in every run and is what makes the stability
measurement below well defined; on the real datasets this is exactly the
resampling protocol the data permits, and applying it to the synthetic datasets
as well keeps all four benchmarks comparable.

Every reported quantity is the mean over replicates, and every error bar is the
$t$-based 95\% confidence interval on that mean. With five replicates these
intervals are wide, and they are meant to be read as such: a difference between
two methods whose intervals overlap is not evidence of a real difference, and
several of the gaps reported below fall into that category. This is a deliberate
correction to the single-run design of the earlier draft, where apparent
rankings could not be distinguished from seed noise.

Replication also permits a direct measurement of the property this literature
calls stability but rarely quantifies. For each method we compute the mean
pairwise Jaccard index between its own selected sets across the replicate seeds,
\begin{equation}
\mathrm{Stab}(M) = \binom{R}{2}^{-1} \sum_{r < r'} \frac{|S_M^{(r)} \cap S_M^{(r')}|}{|S_M^{(r)} \cup S_M^{(r')}|},
\end{equation}
over $R$ replicates. A value of 1 means the method returns the same feature set
every time; a value near 0 means the returned signature is essentially a fresh
draw on each run, which would undermine any biological interpretation of it.

\subsubsection{Performance Metrics}

For predictive performance we report $R^2$, RMSE and MAE for regression; accuracy, F1, precision,
recall and AUC-ROC for binary classification; and accuracy, macro F1 and weighted F1 for multiclass.
Feature recovery on the synthetic datasets is measured through TPR, FDR, precision and selection F1.
Computational efficiency is the wall-clock feature selection runtime in minutes. Sparsity is the
number of features selected, $|S|$, and the sparsity ratio $|S|/p$. Cross-method agreement is
measured by the Jaccard similarity $|S_i \cap S_j| / |S_i \cup S_j|$ and the overlap coefficient
$|S_i \cap S_j| / \min(|S_i|, |S_j|)$.

All of these are reported as a mean over the replicate seeds with a confidence interval, per
Section~\ref{sec:replication}. This evaluation framework enables assessment of the central trade-offs
in feature selection: predictive performance against sparsity, sensitivity against specificity in
feature recovery, stability against flexibility, and computational cost against selection quality.

\subsection{Reproducibility}

All methods, data generators, the experiment driver, and the figure and table generation code are
available at \url{https://github.com/KashfiR/SPARCS}. Every result reported below is produced by
\texttt{experiments.py}, which writes a machine-readable record of the selected feature sets,
metrics, runtimes, per-replicate values, and the exact parameter configuration used;
\texttt{figures.py} rebuilds every figure and Table~\ref{tab:main} from that record, so no reported
number is transcribed by hand. Each replicate is checkpointed individually and is identified by a
fingerprint of the settings that produced it, so a run can be spread over several sessions and
additional seeds can be added later without recomputing the existing ones. The synthetic datasets are
generated from fixed seeds and require no download.

\section{Results and Discussion}

\begin{table}[H]
\centering
\caption{Method performance across the four benchmark datasets.}
\label{tab:main}
\small
\IfFileExists{table1.tex}{\input{table1.tex}}{\fbox{\parbox{0.8\textwidth}{\centering\ttfamily
table1.tex not found. Run experiments.py then figures.py.}}}
\end{table}

\subsection{Synthetic Dataset Performance and Feature Recovery}

\subsubsection{Moderately Nonlinear Data (s1)}

On the synthetic dataset with 20\% nonlinearity, SPARCS achieved the highest predictive performance
and the highest True Positive Rate of the four methods, recovering the largest share of truly
informative features. All methods nonetheless exhibited high False Discovery Rates, reflecting the
difficulty of isolating truly informative features in high-dimensional settings. CPSS-Spearman
showed the lowest FDR by selecting very few features in total, but this extreme conservatism came at
the cost of missing the overwhelming majority of true features. SPARCS, ASGDR and Spearman-ISIS
selected at or near the feature budget with similar FDRs, indicating comparable specificity at
matched selection size.

\resultfigure{figures/fig1_feature_recovery.png}{0.95}{Feature recovery (TPR and FDR) on the two
synthetic datasets.}

Figure~\ref{tab:main} and the recovery panels together illustrate that SPARCS and ASGDR achieved the
best balance between TPR and FDR. CPSS-Spearman's very low TPR highlights the trade-off between
stability and sensitivity: while CPSS provides formal error control, it may be overly conservative
for exploratory analysis.

\subsubsection{Highly Nonlinear Data (s2)}

The second synthetic dataset is intentionally more difficult: 80\% nonlinear informative features,
twice as many interaction terms, higher noise, and a lower signal-to-noise ratio. This design tests
method robustness under severe complexity.

Performance patterns revealed a striking finding. ASGDR improved as complexity increased, while all
screening-based methods deteriorated; SPARCS declined and Spearman-ISIS collapsed, performing worst
despite being competitive on the easier dataset. CPSS-Spearman maintained moderate performance.

This counterintuitive result, where adaptive regularization improves while screening methods fail,
reveals that when nonlinearity, interactions, and noise become pervasive, linear methods with strong
regularization may be more robust than sophisticated nonlinear screening. ASGDR's $L_1$ and $L_2$
penalties enforce sparsity and stability simultaneously, preventing overfitting to spurious nonlinear
patterns. In contrast, screening methods attempt to detect marginal and conditional associations that
become increasingly unreliable as noise dominates signal. The interaction terms further confound
marginal screening, since features important only through an interaction appear marginally weak and
are discarded.

Feature recovery deteriorated across the board in this scenario. ASGDR achieved the highest TPR,
CPSS-Spearman again selected very few features with minimal recovery, and False Discovery Rates
increased for every method. The comparable TPR between SPARCS and Spearman-ISIS suggests that
RDC-based nonlinear screening was insufficient to overcome the dominance of highly nonlinear effects,
though SPARCS maintained better predictive performance than Spearman-ISIS.

\subsubsection{Computational Efficiency on Synthetic Data}

Runtime analysis reveals stark differences in computational efficiency. ASGD is consistently fastest,
leveraging efficient SGD optimization. Spearman-ISIS is also rapid through fast marginal screening.
SPARCS requires moderate time due to RDC computation, permutation testing, and batch CPSS.
CPSS-Spearman is slowest, requiring 60 applications of the base selector (30 pairs $\times$ 2
halves). Despite being substantially slower than Spearman-ISIS, SPARCS delivers better predictive
performance and feature recovery on the moderately nonlinear dataset, a favorable trade-off between
computation and selection quality for discovery work.

\subsection{Real Cancer Dataset Performance}

\subsubsection{Binary Classification: Liver Cancer}

For binary classification of hepatocellular carcinoma versus normal liver tissue (357 samples,
22{,}278 probes), all four methods achieved excellent predictive performance, demonstrating the
viability of the selected feature sets for clinical classification. SPARCS, Spearman-ISIS and
CPSS-Spearman achieved effectively identical accuracy, F1 and recall, with near-identical precision;
Spearman-ISIS achieved a marginally higher AUC. ASGDC showed slightly lower but still strong
performance.

The similar performance across methods suggests that for this well-separated binary classification
problem, multiple feature subsets can achieve high accuracy. Three of the four methods converged on
the same selection size, with CPSS-Spearman slightly larger, indicating broad agreement on the
appropriate sparsity level for this dataset.

\resultfigure{figures/fig2_performance_vs_runtime.png}{0.95}{Test-set performance against feature
selection runtime.}

The performance-runtime panel shows ASGDC achieving comparable performance at a small fraction of the
cost, while CPSS-Spearman is by far the slowest, with SPARCS and Spearman-ISIS between the two
extremes. For production deployment ASGDC's speed advantage is substantial; for research requiring
rigorous error control, CPSS-Spearman's runtime may be acceptable.

\subsubsection{Multiclass Classification: Leukemia Subtype Discrimination}

The 7-class leukemia subtype task (281 samples, 22{,}284 probes) separated the methods more clearly
than the binary case. ASGDC achieved the highest accuracy and macro F1, substantially outperforming
the others. SPARCS achieved the second-best macro F1, with Spearman-ISIS and CPSS-Spearman
performing similarly but slightly worse.

Notably, SPARCS terminated on its own stopping criterion well below the feature budget, selecting
roughly half as many features as the other methods while remaining competitive. This sparsity is a
consequence of the adaptive threshold and stopping rule of Section 2.2.1: selection ends when no
candidate clears both the stability and RDC magnitude thresholds, rather than continuing until a
fixed target is met.

ASGDC's stronger performance while selecting more features than SPARCS indicates that adaptive
regularization may be particularly effective for multiclass problems, since the $L_2$ penalty
stabilizes coefficient estimation across multiple classes, though it is less useful where high
sparsity and human interpretability are the priority. ASGDC was also the fastest method, making it
attractive for multiclass applications where computational efficiency is critical.

\subsection{Method Stability and Feature Set Overlap}

\subsubsection{Cross-Method Agreement Patterns}

\resultfigure{figures/fig3_jaccard_heatmaps.png}{0.95}{Jaccard similarity heatmaps quantifying the
overlap between every pair of methods.}

\paragraph{Spearman-ISIS and CPSS-Spearman show the highest agreement on the real datasets} but
markedly lower agreement on the synthetic ones. Strong agreement on real data is expected, since
CPSS-Spearman uses Spearman-ISIS as its base selector and retains the features it selects most
frequently; the high overlap confirms that CPSS is identifying the stable core of ISIS's selections.
The weaker agreement on synthetic data suggests either that the synthetic response lacks the shared
structure that unifies the two methods on real data, or that the synthetic dependence patterns are
simpler than their biological counterparts. Either way, the divergence is a caution against reading
synthetic feature-recovery numbers as a proxy for real-data behavior.

\paragraph{ASGD shows minimal overlap with screening-based methods on the real datasets}, with
substantially higher overlap on synthetic data. This low overlap indicates that adaptive
regularization and iterative screening identify fundamentally different feature subsets, and again
suggests that the synthetic data may not be representative. ASGD's linear optimization prioritizes
features with strong marginal and conditional linear effects, while screening methods, and especially
SPARCS with RDC, can detect features with weaker linear but stronger nonlinear associations.

\paragraph{SPARCS shows moderate agreement with both paradigms.} Its intermediate position reflects
its hybrid nature: it incorporates elements of both screening (ISIS-style iteration) and stability
(CPSS) while adding nonlinear detection (RDC). The moderate overlap suggests SPARCS identifies a
partially distinct feature set, potentially capturing nonlinear features missed by purely linear or
purely marginal methods.

\paragraph{Overlap decreases with increasing nonlinearity.} Comparing the 20\% and 80\% nonlinear
synthetic datasets, Jaccard indices fall for most method pairs. Methods diverge more when nonlinear
effects dominate, as different approaches prioritize different aspects of the feature-response
relationship.

\subsubsection{Detailed Feature Intersection Analysis}

\resultfigure{figures/fig4_intersections.png}{0.95}{Number of features shared by each pair of methods
across the four datasets.}

The SPARCS and Spearman-ISIS intersection is moderate throughout, indicating that although the two
share a rank-based screening foundation, the additional SPARCS components (RDC, permutation testing,
batch CPSS) shift the selected set considerably: partial methodological alignment with distinct
selection priorities.

The SPARCS and CPSS-Spearman intersection is highly dataset-dependent. On the synthetic datasets the
intersection is limited by CPSS's extreme conservatism, so it accounts for nearly all of CPSS's
selections but only a small fraction of SPARCS's. On the real datasets overlap increases
substantially, suggesting better agreement when the signal-to-noise ratio is higher. The overlap is
asymmetric because the two methods select sets of very different size.

Multiclass classification showed the most distinct pattern, with SPARCS selecting a small set that is
largely contained in the larger sets chosen by the two screening methods. This high containment
suggests SPARCS identified a core subset of features that the other methods also recognized as
important, achieving comparable accuracy with superior parsimony. The very low SPARCS-ASGDC
intersection points to fundamentally distinct feature priorities between integrated screening and
adaptive regularization.

\subsubsection{Selection Stability Under Reseeding}

\resultfigure{figures/fig6_selection_stability.png}{0.8}{Selection stability: the mean pairwise
Jaccard index between a method's own selected sets across replicate seeds. Higher is more
reproducible.}

Cross-method overlap answers whether two methods agree with each other; it does not answer whether a
single method agrees with itself. The replicate design lets us measure the latter directly, and it
is the more consequential quantity for biomarker work, since a signature that changes when only the
split changes cannot carry a biological interpretation regardless of how well it predicts. The
stability values should be read alongside the confidence intervals in
Table~\ref{tab:main}: a method can be both accurate on average and unstable in what it returns, and
those are separate failure modes with separate consequences.

\subsubsection{Implications for Method Selection Stability}

The stability analysis reveals a critical trade-off. \textbf{CPSS-Spearman provides maximal
stability} by construction, having the highest self-agreement across subsamples, but
\textbf{sacrifices sensitivity}: very few features are selected and TPR on synthetic data is low.
\textbf{SPARCS balances stability and performance}, incorporating CPSS-based stability assessment
while using adaptive thresholds and batch addition to maintain reasonable sensitivity. Its moderate
Jaccard indices place it neither at the maximally stable end (like CPSS) nor at the maximally
flexible end (like pure screening), a middle ground that may be well suited to exploratory biomarker
discovery.

The low cross-paradigm agreement between screening and regularization raises important questions
about feature selection consensus. When different valid methods select largely non-overlapping
feature sets yet achieve similar predictive performance, it suggests either that multiple equivalent
feature sets exist (redundancy in the feature space) or that different methods capture complementary
aspects of the feature-response relationship. The consistent performance of diverse feature sets in
binary classification supports the redundancy hypothesis for that dataset, while the performance
differences in multiclass classification suggest complementary information.

\subsection{Sparsity and Computational Trade-offs}

\resultfigure{figures/fig5_sparsity_runtime.png}{0.95}{Selection sparsity and runtime across methods
and datasets. Runtime is on a logarithmic scale.}

CPSS-Spearman exhibits extreme variability in sparsity between the synthetic and real datasets. This
instability stems from its selection frequency threshold: when the signal is weak, few features
exceed $\tau$, whereas when the signal is strong many features appear stable. This dataset-dependent
behavior limits CPSS's utility when a desired sparsity level is predetermined.

SPARCS shows more consistent sparsity, with the notable exception of multiclass classification where
the adaptive threshold identified a natural stopping point well below the budget. This data-driven
sparsity is desirable: SPARCS adapts to the structure of each dataset rather than imposing a fixed
feature count. Spearman-ISIS and ASGD select up to their budget by construction, giving predictable
sparsity but potentially missing dataset-specific optimal sizes.

Runtime differences span orders of magnitude, with ASGD fastest, then Spearman-ISIS, then SPARCS,
then CPSS-Spearman. The gap between the fastest and slowest methods is close to three orders of
magnitude and has practical consequences: for routine clinical deployment ASGD's near-instantaneous
selection is highly attractive, while for one-time biomarker discovery studies SPARCS's runtime is
acceptable given its stronger feature recovery and balanced performance.

\subsection{SPARCS: Strengths and Limitations}

SPARCS integrates rank-based SIS, RDC residual screening, permutation testing, and batch CPSS. This
multilayered design improved feature recovery in moderately complex settings, achieving the highest
TPR alongside competitive predictive performance. It balances sensitivity and specificity better than
CPSS, which is overly conservative, by using adaptive thresholding to set sparsity from the data
rather than from a fixed target. Its runtime is acceptable for discovery studies where improved
recovery outweighs the extra compute.

SPARCS struggles under extreme nonlinearity: on the 80\% nonlinear dataset it produced negative $R^2$
and low recovery, whereas ASGDR improved. RDC can miss complex interactions or threshold effects, and
its bandwidth $s$ trades sensitivity to high-frequency dependence against the level of the null. The
computational overhead of RDC, permutations and CPSS is only justified when it finds signal that
simpler methods miss; on easy, well-separated problems simpler methods suffice. Multiple
hyperparameters ($M_{\mathrm{pre}}$, $M_{\mathrm{rdc}}$, $n_{\mathrm{perm}}$, $\alpha$, $\gamma$,
$k_{\mathrm{add}}$, the RDC threshold) affect performance, and theoretical guarantees remain limited
compared with CPSS. SPARCS can also be overkill for straightforward binary classification.

\section{Discussion}

\subsection{Main Paradigm Takeaways}

Adaptive regularization (ASGD) is the fastest method and the most robust in highly complex and noisy
settings. It favors features with strong linear contributions and has low overlap with screening
methods. Greedy iterative screening (Spearman-ISIS) provides an excellent speed-accuracy balance and
an interpretable selection path, but lacks any stability assessment. Stability selection
(CPSS-Spearman) guarantees maximal stability with theoretical backing, but can be too conservative
for discovery in weak-signal settings. SPARCS, integrating screening and stability, occupies a middle
ground: better recovery than pure screening at moderate complexity, more sensitive than CPSS, but not
always superior under pervasive nonlinearity. SPARCS is most useful for exploratory biomarker
discovery in noisy or nonlinear problems where sensitivity matters and moderate runtimes are
acceptable.

Low cross-paradigm overlap implies that multiple equally predictive but different feature sets exist,
reflecting either redundancy or complementary signal. For robust biomarkers, consensus or ensemble
approaches, such as the intersection of SPARCS and ASGDR, or their union filtered by CPSS stability,
would prioritize features with higher reproducibility.

Unlike CPSS-Spearman, which achieves low FDR through extreme conservatism, SPARCS maintains moderate
FDR while selecting substantially more features. This balance matters for exploratory biomarker
discovery, where missing true features is more costly than including false positives that can be
filtered in subsequent validation pipelines.

\subsection{Limitations}

Three limitations bound the conclusions. First, the synthetic datasets differ markedly from the real
omics datasets in the cross-method overlap they induce, so synthetic feature-recovery numbers should
be read as a controlled diagnostic rather than a prediction of real-data behavior. Second, five replicates is enough to expose seed sensitivity but not enough to
estimate it precisely; the confidence intervals are correspondingly wide, and differences smaller
than those intervals should not be read as rankings. A larger replicate budget, or a paired test
across seeds, would sharpen the comparisons that currently overlap. Third, the comparison covers two real datasets from a single curated repository and
one downstream predictor, so generalization to other assay types, disease areas, and model classes is
untested.

\subsection{Future Work}

Broader benchmarking across more datasets and against tree-based, deep and Bayesian selectors is
needed. In particular, classification with more realistic synthetic omics data, specifically using
Model-X knockoffs, should be tested, and extensive hyperparameter tuning of SPARCS's parameters will
be pursued \citep{zhou2025}. Alternative dependence measures such as the Hilbert-Schmidt independence
criterion and the computationally intensive distance correlation, artificial neural network
screening, adaptive hyperparameter selection, and theoretical analysis of SPARCS for Type I error
control and sure-screening conditions will also be explored \citep{gretton2005,hou2022}. Integrating
uncertainty quantification such as bootstrap confidence intervals or Bayesian posteriors would
improve interpretability and clinical translation \citep{justus2024}. Increasing the replicate budget beyond five seeds, and adding paired
significance testing across seeds, is the most immediate priority for the comparisons whose
intervals currently overlap.

\section{Conclusions}

This paper contributes an empirical, systematic comparison of four paradigms (the novel SPARCS,
Spearman-ISIS, CPSS-Spearman, and ASGD) across synthetic and real cancer genomics datasets to
highlight trade-offs in feature recovery, stability, runtime, and parsimony. Feature selection
remains a central challenge in genomic medicine, as robust biomarker identification is essential for
precision diagnostics and targeted therapies. Our results show that no single method is universally
optimal; performance depends on dataset characteristics, computational constraints, and application
goals, with all current approaches achieving only moderate success in complex, nonlinear settings.
Future advances are likely to arise from integrating biological prior knowledge, developing methods
tailored to nonlinear and interaction-rich signals including deep learning approaches, and using
ensemble strategies to capture complementary information. Within this landscape, SPARCS represents a
step toward integrative feature selection, enabling more reliable and interpretable biomarker
discovery that can improve clinical decision-making, enhance diagnostic and prognostic prediction,
accelerate drug development, and reduce healthcare costs through smaller, more precise molecular
panels.

\bibliographystyle{unsrtnat}
\begin{thebibliography}{99}
\small
\bibitem{hastie2009} Hastie, T., Tibshirani, R., \& Friedman, J. (2009). \textit{The Elements of Statistical Learning: Data Mining, Inference, and Prediction} (2nd ed.). Springer.
\bibitem{saeys2007} Saeys, Y., Inza, I., \& Larra\~naga, P. (2007). A review of feature selection techniques in bioinformatics. \textit{Bioinformatics}, 23(19), 2507--2517.
\bibitem{guyon2003} Guyon, I., \& Elisseeff, A. (2003). An introduction to variable and feature selection. \textit{Journal of Machine Learning Research}, 3, 1157--1182.
\bibitem{buhlmann2011} B\"uhlmann, P., \& van de Geer, S. (2011). \textit{Statistics for High-Dimensional Data: Methods, Theory and Applications}. Springer.
\bibitem{barabasi2011} Barab\'asi, A. L., Gulbahce, N., \& Loscalzo, J. (2011). Network medicine: a network-based approach to human disease. \textit{Nature Reviews Genetics}, 12(1), 56--68.
\bibitem{zhao2014} Zhao, S. D., Cai, T. T., \& Li, H. (2014). Direct estimation of differential networks. \textit{Biometrika}, 101(2), 253--268.
\bibitem{liu2009} Liu, H., Lafferty, J., \& Wasserman, L. (2009). The nonparanormal: semiparametric estimation of high dimensional undirected graphs. \textit{Journal of Machine Learning Research}, 10, 2295--2328.
\bibitem{anders2010} Anders, S., \& Huber, W. (2010). Differential expression analysis for sequence count data. \textit{Genome Biology}, 11(10), R106.
\bibitem{chandrashekar2014} Chandrashekar, G., \& Sahin, F. (2014). A survey on feature selection methods. \textit{Computers \& Electrical Engineering}, 40(1), 16--28.
\bibitem{li2017} Li, J., Cheng, K., Wang, S., Morstatter, F., Trevino, R. P., Tang, J., \& Liu, H. (2017). Feature selection: a data perspective. \textit{ACM Computing Surveys}, 50(6), 1--45.
\bibitem{tibshirani1996} Tibshirani, R. (1996). Regression shrinkage and selection via the lasso. \textit{Journal of the Royal Statistical Society: Series B}, 58(1), 267--288.
\bibitem{zou2005} Zou, H., \& Hastie, T. (2005). Regularization and variable selection via the elastic net. \textit{Journal of the Royal Statistical Society: Series B}, 67(2), 301--320.
\bibitem{zou2006} Zou, H. (2006). The adaptive lasso and its oracle properties. \textit{Journal of the American Statistical Association}, 101(476), 1418--1429.
\bibitem{fan2008} Fan, J., \& Lv, J. (2008). Sure independence screening for ultrahigh dimensional feature space. \textit{Journal of the Royal Statistical Society: Series B}, 70(5), 849--911.
\bibitem{fan2009} Fan, J., Samworth, R., \& Wu, Y. (2009). Ultrahigh dimensional feature selection: beyond the linear model. \textit{Journal of Machine Learning Research}, 10, 2013--2038.
\bibitem{fan2010} Fan, J., \& Song, R. (2010). Sure independence screening in generalized linear models with NP-dimensionality. \textit{The Annals of Statistics}, 38(6), 3567--3604.
\bibitem{li2012} Li, G., Peng, H., Zhang, J., \& Zhu, L. (2012). Robust rank correlation based screening. \textit{The Annals of Statistics}, 40(3), 1846--1877.
\bibitem{wang2013} Wang, X., Jiang, Y., Huang, M., \& Zhang, H. (2013). Robust variable selection with exponential squared loss. \textit{Journal of the American Statistical Association}, 108(502), 632--643.
\bibitem{croux2010} Croux, C., \& Dehon, C. (2010). Influence functions of the Spearman and Kendall correlation measures. \textit{Statistical Methods \& Applications}, 19(4), 497--515.
\bibitem{meinshausen2010} Meinshausen, N., \& B\"uhlmann, P. (2010). Stability selection. \textit{Journal of the Royal Statistical Society: Series B}, 72(4), 417--473.
\bibitem{nogueira2018} Nogueira, S., Sechidis, K., \& Brown, G. (2018). On the stability of feature selection algorithms. \textit{Journal of Machine Learning Research}, 18(1), 6345--6398.
\bibitem{shah2013} Shah, R. D., \& Samworth, R. J. (2013). Variable selection with error control: another look at stability selection. \textit{Journal of the Royal Statistical Society: Series B}, 75(1), 55--80.
\bibitem{eindor2005} Ein-Dor, L., Kela, I., Getz, G., Givol, D., \& Domany, E. (2005). Outcome signature genes in breast cancer: is there a unique set? \textit{Bioinformatics}, 21(2), 171--178.
\bibitem{michiels2005} Michiels, S., Koscielny, S., \& Hill, C. (2005). Prediction of cancer outcome with microarrays: a multiple random validation strategy. \textit{The Lancet}, 365(9458), 488--492.
\bibitem{hofner2015} Hofner, B., Boccuto, L., \& G\"oker, M. (2015). Controlling false discoveries in high-dimensional situations: boosting with stability selection. \textit{BMC Bioinformatics}, 16(1), 144.
\bibitem{lopezpaz2013} Lopez-Paz, D., Hennig, P., \& Sch\"olkopf, B. (2013). The randomized dependence coefficient. \textit{Advances in Neural Information Processing Systems}, 26, 1--9.
\bibitem{ernst2004} Ernst, M. D. (2004). Permutation methods: a basis for exact inference. \textit{Statistical Science}, 19(4), 676--685.
\bibitem{salvador2004} Salvador, S., \& Chan, P. (2004). Determining the number of clusters/segments in hierarchical clustering/segmentation algorithms. In \textit{Proceedings of the 16th IEEE International Conference on Tools with Artificial Intelligence} (pp. 576--584).
\bibitem{feltes2019} Feltes, B. C., Chandelier, E. B., Grisci, B. I., \& Dorn, M. (2019). CuMiDa: an extensively curated microarray database for benchmarking and testing of machine learning approaches in cancer research. \textit{Journal of Computational Biology}, 26(4), 376--386.
\bibitem{ke2017} Ke, G., Meng, Q., Finley, T., Wang, T., Chen, W., Ma, W., Ye, Q., \& Liu, T. Y. (2017). LightGBM: a highly efficient gradient boosting decision tree. \textit{Advances in Neural Information Processing Systems}, 30, 3146--3154.
\bibitem{zhou2025} Zhou, J., \& Scherr, S. (2025). Model-X knockoffs in the replication crisis era: reducing false discoveries and researcher bias in social science research. \textit{Social Sciences \& Humanities Open}, 11, 101380.
\bibitem{gretton2005} Gretton, A., Bousquet, O., Smola, A., \& Sch\"olkopf, B. (2005). Measuring statistical dependence with Hilbert-Schmidt norms. \textit{Lecture Notes in Computer Science}, 63--77.
\bibitem{hou2022} Hou, J., Ye, X., Feng, W., Zhang, Q., Han, Y., Liu, Y., Li, Y., \& Wei, Y. (2022). Distance correlation application to gene co-expression network analysis. \textit{BMC Bioinformatics}, 23(1), 81.
\bibitem{justus2024} Justus, V. L., Rodrigues, V. B., \& Sousa, A. R. S. (2024). Bootstrap confidence intervals: a comparative simulation study. \textit{arXiv:2404.12967}.
\end{thebibliography}

\end{document}
